\subsection{Different ideals and a split quadratic dimension field}
\label{sec:round3-different}
The remaining dense quadratic profile requires a refinement that also controls
possible larger signed dimension fields.

\begin{theorem}[Different-ideal obstruction]\label{thm:signed-different}
Let $K$ be the dimension-character field of the commutative signed algebra
obtained by sphericalization of a commutative self-dual fusion ring.
Its global dimension $\Delta$ belongs to the different
$\mathfrak D_{K/\mathbb Q}$.
\end{theorem}
\begin{proof}
Write $d_i$ for the signed basis dimensions, and let $e_K\in K\otimes K$
be the separability idempotent. Orthogonality of distinct Galois-conjugate
characters gives
\[
 (\Delta\otimes1)e_K=\sum_i d_i\otimes d_i
                 \in\mathcal O_K\otimes\mathcal O_K.
\]
For $a$ in the codifferent, contract the second factor by
$x\mapsto\operatorname{Tr}(ax)$. The right side is integral, and the left
side is $\Delta a$. Thus $\Delta\mathfrak D^{-1}\subseteq\mathcal O_K$,
which is the assertion.
\end{proof}

For a local extension with ramification index $e$ and different exponent $d$,
\[
 \operatorname{Tr}(\mathfrak p^k)
       =p^{\lfloor(k+d)/e\rfloor}\mathbb Z_p.
\]
In particular, every tamely ramified odd prime divides the trace of the
different. A cyclic cubic extension of conductor $9$ has $e=3,d=4$ and
contributes $9$. We use these standard local trace formulas and the
conductor--discriminant formula; see Serre, \emph{Local Fields},
Chapters III and VI, and Milne, \emph{Class Field Theory}, Chapter V, Theorem~3.27.\footnote{The latter reference is available at
\url{https://www.jmilne.org/math/CourseNotes/CFT.pdf}.}

\begin{proposition}\label{prop:split-41-profile}
A commutative self-dual rank-six fusion ring with formal codegrees
\[
 9,9,9,9,\qquad 3690\pm576\sqrt{41},
\]
Casimir trace $7416$, and reduction
$\mathbb F_2^4\times\mathbb F_2[\epsilon]/(\epsilon^2)$ satisfies the following
necessary condition for categorifiability: its signed Casimir trace is $7416$.
Consequently every signed coefficient has its original absolute value.
\end{proposition}
\begin{proof}
Put $E=\mathbb Q(\sqrt{41})$. The dimension-field containment lemma gives
$E\subseteq K$, $\Delta/9\in\mathcal O_K$, and
\[
 \frac{\Delta}{18\sqrt{41}}\in\mathcal O_K,
 \qquad 738\mid\operatorname{Tr}_{K/\mathbb Q}(\Delta),
 \qquad \operatorname{Tr}_{K/\mathbb Q}(\Delta)\le7416.
\]
Indeed $3690+576\sqrt{41}=18\sqrt{41}(32+5\sqrt{41})$, and the last factor
is an integral unit of norm $-1$. Cyclotomicity makes $K/\mathbb Q$
abelian Galois, of degree $2$, $4$, or $6$.

If $[K:\mathbb Q]$ is $4$ or $6$, then $2$ splits completely in $K$.
The four etale rank-one blocks of the completed signed order have generic
fibre $\mathbb Q_2$; its only other block has rank two. A Galois number
field using a completion of degree two would require at least two such
local components, whereas only one rank-two block is available. The same
argument excludes ramification at $2$. It allows rational idempotents
inside the non-etale block and does not assume that block is maximal.

Theorem~\ref{thm:signed-different} and the trace bound restrict ramified
primes outside $\{3,41\}$ to $5$ and $7$, and forbid their simultaneous
occurrence, which would require trace at least $35\cdot738$.
The complete list of possible larger fields follows from elementary
Dirichlet-character arithmetic:
\begin{enumerate}
\item A biquadratic field $E\mathbb Q(\sqrt m)$ requires $m>1$ squarefree,
positive, prime to $41$, supported on $\{3,5,7\}$, and $m\equiv1\pmod8$.
The only such product is $105$, already forbidden by its primes $5$ and $7$.
\item A cyclic quartic character with square the quadratic character of
conductor $41$ is a quartic character modulo $41$ times a quadratic
character of positive odd discriminant $m$. Since $2^{10}\equiv-1\pmod{41}$,
splitting at $2$ requires $m\equiv5\pmod8$. The permitted products leave
exactly $m=5,21$. The conductors are $205,861$, and the possible dimension
traces are respectively $3690,7380$ and $5166$.
\item A degree-six field is $E L$ with $L$ cyclic cubic. Its conductor is
$7$, $9$, or $63$: neither $5$ nor $41$ admits a cubic character, and a
nontrivial cubic character at $3$ has conductor $9$. The prime $2$ is inert
in the conductor-$7$ and conductor-$9$ cubic fields. Only conductor $63$
remains, with dimension trace $5166$.
\end{enumerate}
There are no additional prime-power conductors: characters of order dividing
$4$ or $6$ are trivial on the higher prime-power unit groups at primes other
than $2$ or $3$; the relevant quadratic odd characters have squarefree conductor.

In the sextic case, coprimality of the discriminants gives
$\mathcal O_{EL}=\mathcal O_E\otimes\mathcal O_L$ and a product of differents.
Write
\[
 \Delta=18\sqrt{41}\,\frac{a+b\sqrt{41}}2,
 \qquad a,b\in\mathcal O_L,\quad a\equiv b\pmod2.
\]
Positivity at both embeddings of $E$ over each real embedding of $L$ makes
$b$ totally positive, with $\operatorname{Tr}_{L/\mathbb Q}(b)=7$.
At $7$, the cubic different is $\mathfrak p_7^2$, and $18$ is a unit.
The different theorem therefore forces $b\in\mathfrak p_7^2$. Hence
\[
 49\le\operatorname{Norm}(b)\le(7/3)^3<49,
\]
a contradiction.

The cyclic quartic cases require a short exact lattice enumeration. For
$M=5,21$, set
\[
 \alpha^2=\frac{M(41+5\sqrt{41})}{2},\qquad
 \beta^2=\frac{M(41-5\sqrt{41})}{2},\qquad
 \alpha\beta=2M\sqrt{41}.
\]
The defining polynomial is $X^4-41MX^2+164M^2$, and an integral basis is
\[
 1,\quad\frac{1+\sqrt{41}}2,\quad
 \frac{1+\sqrt{41}+2\alpha}4,\quad
 \frac{1-\sqrt{41}+2\beta}4.
\]
The checker verifies its integral multiplication, cyclic Galois action,
complete splitting at $2$, and discriminant $41^3M^2$. The preceding field
list and the conductor--discriminant formula prove maximality.
For $\theta=\Delta/(18\sqrt{41})=\sum_jz_jb_j$ and trace-dual basis $b_j^*$,
\[
 z_j=\operatorname{Tr}\left(\Delta\frac{b_j^*}{18\sqrt{41}}\right).
\]
At fixed positive trace $T$, this lies between $T$ times the smallest and
largest conjugates of $b_j^*/(18\sqrt{41})$. Exact algebraic sign comparisons
certify a finite integer box for every $z_j$. Imposing the trace and total
positivity leaves $28,216,16$ candidates at traces $3690,7380,5166$,
respectively. Every candidate fails a recorded $d$-number divisibility.
The supplementary certificate contains all $260$ candidates and the exact
coefficient witnessing each failure.

It remains to consider $K=E$. Write
$\Delta=18(41y+x\sqrt{41})/2$. Then $x\equiv y\pmod2$, $1\le y\le10$,
$x^2<41y^2$, and
\[
 41y^2-x^2\mid164y^2.
\]
The last condition is the quadratic $d$-number test. The four complementary
codegrees form either four rational integers, two rational integers and a
quadratic pair, or two quadratic pairs in $E$. Rational codegrees divide
the integral content of $\Delta$. A quadratic pair has polynomial
$X^2-SX+P$, with $S^2-4P=41r^2$, $P\mid\operatorname{Norm}(\Delta)$,
$P\mid S^2$, both roots greater than one, and both dividing $\Delta$
integrally. The total trace bounds $S$ and $r$. The reciprocal identity
$\sum_f1/f=1$ completes the finite test. Exactly seven dimension candidates
reach this stage; none admits a quadratic complementary pair. The only
complete spectra have $y=10$, $x=\pm64$, and four codegrees $9$.
Their total trace is $7416$, proving the proposition.
\end{proof}

The full tensor of \texttt{R6\_8\_6526} satisfies these hypotheses.
Exact saturated signed associativity leaves one table up to the certified
sign gauges, equivalent to the original tensor. Every non-FP original
character orbit fails the arithmetic dimension-character tests. Thus a
categorification would have a pseudo-unitary Galois conjugate, which the
existing global-codegree moment certificate excludes. The resulting
\emph{general} exclusion is implemented in
\filename{software/round3_split_arithmetic.py}, with canonical certificate
\filename{research/round3/arithmetic/R6_8_6526_general.json}.
